\section{Link layer}

\begin{knBox}
    {Link layer services}
    \begin{itemize}
        \item \textbf{Framing}: encapsulating network-layer datagrams into frames, adding headers and trailers
        \item \textbf{Link access}: controlling access to the shared physical medium (e.g., CSMA/CD for Ethernet)
        \item \textbf{Reliable delivery}: retransmission, error detection and correction to ensure frames are delivered without errors (e.g., CRC)
    \end{itemize}
\end{knBox}

\begin{knBox}
    {Network Interface Card (NIC)}
    The link layer is implemented in the NIC.
\end{knBox}

\begin{definition}
    {MAC Addressing}
    \begin{itemize}
        \item Used to get frame from one interface to another physically-connected inferface under the same \textbf{local area network}
        \item Burned into NIC by manufacturer. Portable.
        \item Unique in network
        \item 48-bit
        \item Also called the \textit{physical} or \textit{link-layer} address
    \end{itemize}
    \begin{itemize}
        \item The source and destination MAC addresses of a packet is \textbf{changed} each time \textbf{it is sent} (passing through a router).
    \end{itemize}
\end{definition}

\begin{knBox}
    {MAC Address Assignment}
    \begin{itemize}
        \item Manufacturers purchase blocks of MAC addresses from IEEE
        \item Portable: can move NIC from one computer to another
    \end{itemize}
\end{knBox}

\subsection{Error Detection}

\begin{theorem}
    {Single Bit Parity Check}
    Add 1 more bit to the data, value of which is set such that the total number of 1's in the data + parity bit is even (or odd).
\end{theorem}

\begin{theorem}
    {Two Dimensional Parity Check}
    Arrange data in a matrix. Add $r+c+1$ parity bits, where $r$ is the number of rows and $c$ is the number of columns.

    Additional bits \texttt{(each row, each column, overall)} ensure even (or odd) parity.
\end{theorem}

\begin{theorem}
    {Cyclic Redundancy Check (CRC)}
    Given data $d$, bitstring $b$, CRC bit size $s$, compute CRC bits as follows:
    \begin{enumerate}
        \item Append $s$ zeros to $d$ at end to be $s_2$
        \item Divide $s_2$ by $b$ using XOR binary division (\textit{long divison in base 2 using XOR (01->1) instead of the usual subtract})
        \item Remainder $r$ is CRC bits
    \end{enumerate}
\end{theorem}

\sectionbar{M07-L01}

\subsection{Access control and ARP}

\subsubsection{Access Control Protocols}

\begin{definition}
    {CSMA (Carrier Sense Multiple Access)}
    \begin{itemize}
        \item Listen before sending; transmit only if the medium is idle.
        \item A random-access protocol that matches human conversational conventions.
    \end{itemize}
\end{definition}

\begin{definition}
    {CSMA/CD (Carrier Sense Multiple Access with Collision Detection)}
    \begin{itemize}
        \item While transmitting, adapters listen for collisions.
        \item If a collision is detected, abort transmission immediately to save time.
        \item Behaviour also mirrors human conversations: both speakers stop and resolve who goes first.
    \end{itemize}
\end{definition}

\begin{definition}
    {Binary Exponential Backoff}
    \begin{itemize}
        \item After a collision, choose a random backoff interval before retrying.
        \item The backoff is chosen uniformly from the range $[0,\,2^{n}-1]$, where $n$ is the number of consecutive collisions for that frame.
        \item This adapts waiting time: small contention → short waits; large contention → longer waits.
    \end{itemize}
\end{definition}

\begin{definition}
    {Turn-Based Access Control}
    \begin{itemize}
        \item Senders take turns transmitting; if a sender has nothing to send, it passes its turn.
        \item Can be implemented with either centralized control or decentralized (token-based) control.
    \end{itemize}
\end{definition}

\begin{definition}
    {Centralized Control}
    \begin{itemize}
        \item A single controller decides whose turn it is (polling or explicit request).
        \item Example: Bluetooth master device controls when peripherals may send.
    \end{itemize}
\end{definition}

\begin{definition}
    {Decentralized Control (Token Passing)}
    \begin{itemize}
        \item A token is passed among nodes; only the token holder may transmit.
        \item Complications: token loss or duplicate tokens require recovery (token regeneration/election).
    \end{itemize}
\end{definition}

\subsubsection{Switches}

\begin{definition}
    {Switch}

    \begin{itemize}
        \item Link-layer device that learns the MAC addresses of connected devices and forwards frames only to the intended recipient, reducing unnecessary traffic.
        \item Contains \textbf{forwarding table} with columns \texttt{(MAC address, interface, timestamp)}.
        \item Entries time out after a certain period of inactivity (e.g., 5 minutes).
        \item Never put broadcast address \texttt{(FF:FF:FF:FF:FF:FF)} in table.
    \end{itemize}

    Incoming frame format: \texttt{Source MAC, Destination MAC, Type, Interface} (Assume Type is always IP)

    Handling process involves two steps. Consider \texttt{sMAC, dMAC, IP, I}:

    \begin{enumerate}
        \item \textbf{Source MAC (Learning):}
              \begin{itemize}
                  \item sMAC in table: replace row
                  \item else: add new row
              \end{itemize}
        \item \textbf{Destination MAC (Forwarding):}
              \begin{itemize}
                  \item dMAC \textit{not} in table: send to all interfaces except \texttt{I}
                  \item else for matching row interface \texttt{I\_0}:
                        \begin{itemize}
                            \item \texttt{I == I\_0}: discard frame
                            \item \texttt{I != I\_0}: forward frame to interface \texttt{I\_0}
                        \end{itemize}
              \end{itemize}
    \end{enumerate}

\end{definition}

\subsubsection{Address Resolution Protocol}

Purpose: resolve IP address to MAC address

\begin{definition}
    {Address Resolution Protocol (ARP)}
    \begin{itemize}
        \item Each host and router on LAN has ARP table mapping IP addresses to MAC addresses.
        \item If host wants to send datagram to IP address not in ARP table, it broadcasts ARP query packet asking "Who has IP address \texttt{x.x.x.x}?"
        \item Host with that IP replies with ARP reply packet containing its MAC address.
        \item Receiver caches sender IP-MAC mapping in its ARP table.
        \item After receiving the reply, sender caches this mapping in its ARP table.
    \end{itemize}
\end{definition}

\begin{theorem}
    {Sending datagram as a new host on LAN}
    \begin{itemize}
        \item To send \colorbox{yellow}{\textbf{inside}} LAN, send ARP request for MAC address of \colorbox{yellow}{\textbf{destination}}.
        \item To send \colorbox{pink}{\textbf{outside}} LAN, send ARP request for MAC address of \colorbox{pink}{\textbf{default gateway}}.
    \end{itemize}
    Process:
    \begin{enumerate}
        \item Host A wants to send datagram to Host B (IP address known, MAC address unknown).
        \item A sends ARP request for IP of address \colorbox{yellow}{B (inside LAN)} / \colorbox{pink}{LAN gateway (outside LAN)}.
              \begin{itemize}
                  \item Source IP and MAC: A's
                  \item Destination MAC: FF:FF:FF:FF:FF:FF (broadcast)
                  \item Destination IP: \colorbox{yellow}{B's IP} / \colorbox{pink}{gateway's IP}
              \end{itemize}
        \item Host B (or default gateway) responds with ARP reply.
              \begin{itemize}
                  \item Source IP and MAC: \colorbox{yellow}{B's} / \colorbox{pink}{gateway's}
                  \item Destination IP and MAC: A's
              \end{itemize}
        \item Datagram:
              \begin{itemize}
                  \item Source IP and MAC: A's
                  \item Destination IP: B's
                  \item Destination MAC: \colorbox{yellow}{B's MAC} / \colorbox{pink}{gateway's MAC}
              \end{itemize}
    \end{enumerate}
\end{theorem}

\begin{knBox}
    {ARP Poisoning}

    Done by swapping expected \textbf{MAC addresses} in ARP table to trudy's MAC address (not IP) to gain access to data.

    This can be eased by \textbf{statically binding} MAC addresses of devices to IP addresses.
\end{knBox}

\sectionbar{M07-L02}

\subsection{Dynamic Host Configuration Protocol}

\begin{definition}
    {Dynamic Host Configuration Protocol (DHCP)}
    \begin{itemize}
        \item \textbf{Application layer} protocol for \textbf{dynamically assigning IP addresses} to new hosts on a network.
        \item Runs on top of \textbf{UDP}.
        \item Allows hosts to obtain the needed info when joining a network:
              \begin{enumerate}
                  \item IP address
                  \item Subnet mask
                  \item Router (default gateway) address
                  \item DNS server addresses
              \end{enumerate}
        \item \texttt{DHCP\_OFFER / DHCP\_ACK} also includes \textbf{lease time} for the assigned IP address.\\
              Send new \texttt{DHCP\_REQUEST} when half lease left. If no response, repeat after some time. Send \texttt{DHCP\_RELEASE} as client when done with lease.
    \end{itemize}
\end{definition}

\begin{theorem}
    {New device joining a network}

    Steps:
    \begin{enumerate}
        \item Client $\rightarrow$ \texttt{DHCP\_DISCOVER} $\rightarrow$ Server
              \begin{itemize}
                  \item src: \texttt{0.0.0.0:68}
                  \item dst: \colorbox{yellow}{\texttt{255.255.255.255:67}}
                  \item yiaddr: \texttt{0.0.0.0}
              \end{itemize}
        \item Client $\leftarrow$ \texttt{DHCP\_OFFER} $\leftarrow$ Server
              \begin{itemize}
                  \item src: \colorbox{pink}{\texttt{server\_ip:67}}
                  \item dst: \colorbox{pink}{\texttt{255.255.255.255:68}}
                  \item yiaddr: \colorbox{pink}{\texttt{offered\_ip}}
              \end{itemize}
        \item Client $\rightarrow$ \texttt{DHCP\_REQUEST} $\rightarrow$ Server
              \begin{itemize}
                  \item src: \texttt{0.0.0.0:68}
                  \item dst: \colorbox{yellow}{\texttt{255.255.255.255:67}}
                  \item yiaddr: \texttt{offered\_ip}
              \end{itemize}
        \item Client $\leftarrow$ \texttt{DHCP\_ACK} $\leftarrow$ Server
              \begin{itemize}
                  \item src: \colorbox{pink}{\texttt{server\_ip:67}}
                  \item dst: \colorbox{pink}{\texttt{255.255.255.255:68}}
                  \item yiaddr: \colorbox{pink}{\texttt{offered\_ip}}
              \end{itemize}
    \end{enumerate}

    Note: all have same transaction ID.

    (Discover $\to$ Offer $\to$ Request $\to$ Ack)
\end{theorem}

\sectionbar{M07-L03}

\subsection{Physical transmission}

Cables

Modulation
AM FM PM

\begin{definition}
    {Virtual Local Area Network (VLAN)}
    \begin{itemize}
        \item Logically separate LANs that share the same physical switch infrastructure.
    \end{itemize}
\end{definition}

\begin{theorem}
    {Switches under switches}
    \begin{itemize}
        \item Full bi-section bandwidth: Hosts within the same switch can communicate at full link bandwidth.
        \item Hosts across different switches share the bandwidth of the inter-switch link.
        \item Throughput limited by bottleneck link, given by \texttt{bandwidth / number of communicating hosts}.
    \end{itemize}
\end{theorem}

\sectionbar{M07-L04}